% =====================================================================
% Figure contract
% 核心结论：中央执行闭环串联任务接口、数据、模型系统、学习、评测与部署开放问题，
%           并由 accountable-state thesis 在 §1、§11 收束全文。
% 图原型：mode A / Skeleton C（中央 hero + 周边 panels）；为 6 panels 与闭环语义重构拓扑。
% 证据层级：hero = CUA 执行闭环；支撑 = §4--§10 六张章节卡；收束 = 底部 summary bar。
% 能不能更少？：不能；删任一章节卡会丢失指定章节导航，删任一 hero 节点会破坏指定闭环。
%
% Step ① 设计文档（Form B / Narrative，复杂综述总图）
% 0 参考：复用 Skeleton C 的中央 hero + 周边卡片节奏，不复刻外部位图。
% 1 领域：Computer-Use Agents 中文综述；技术术语保持英文。
% 2 格式：TikZ standalone，XeLaTeX + fontspec + xeCJK。
% 3 画布：最终 W_ink 约 27 cm；W_print = 24 cm；目标印刷最小 CJK 字号 7.5 pt。
%   设计最小 CJK 字号 >= 8.5 pt；全图自左向右、由章节卡向中央闭环汇聚。
% 4 模块：中央 14.2×8.0 cm hero；左右各 3 张 5.9×2.4 cm panel；底部整宽 summary。
% 5 连线：主闭环实线；R→P feedback；R→H 虚线；学习/人类接管为正交辅助流；
%   panel 到语义节点统一 dotted leader。
% 6 空间：hero 两行主路径；行间 corridor、左右边 rail、底部 rail 分担跨层连线。
% 7 强调：hero 1.4 pt 深色描边 + 独有底色；panels 0.7 pt、六种低饱和 zone tone。
% 8 密度：综述总图复杂档；不造数值/公式/数据可视化，信息 panel 即六章导航与 summary。
% 9 Narrative：左列三卡、中央 hero、右列三卡同高对齐；hero 上两行形成蛇形执行闭环，
%   下层 D/H/L 填满辅助语义；外部空白只保留 0.4--0.6 cm 构造间隙。
% 10 可视化决策：输入无数值，故不嵌入 chart/heatmap；用 loop topology 本身作为视觉 hero。
%
% Pre-flight：P1 设计文档已写；P2 同列/同行按 positioning 构造；P3 全边列于代码；
% P4 N/A（无嵌入数值图）；P5 corridors/rails 已分配；P6 固定尺寸容器；P7 已读 lessons。
% 拓扑适配记录：Skeleton C 的“模型卡 4 panels”改为“执行闭环 + 6 章节卡”；
% 必须新增 R→P、R→H→S、D/V→L→O/P，并使用 node anchors + named waypoints。
% =====================================================================

\documentclass[border=8pt]{standalone}
\usepackage{fontspec}
\usepackage{xeCJK}
\usepackage{xcolor}
\usepackage{tikz}
\IfFontExistsTF{PingFang SC}{
  \setCJKmainfont{PingFang SC}
  \setCJKsansfont{PingFang SC}
}{
  % The same PingFang SC face is present as index 3 of Apple's on-demand TTC.
  \setCJKmainfont[
    Path=/System/Library/AssetsV2/com_apple_MobileAsset_Font8/86ba2c91f017a3749571a82f2c6d890ac7ffb2fb.asset/AssetData/,
    FontIndex=3,
    AutoFakeBold=2
  ]{PingFang.ttc}
  \setCJKsansfont[
    Path=/System/Library/AssetsV2/com_apple_MobileAsset_Font8/86ba2c91f017a3749571a82f2c6d890ac7ffb2fb.asset/AssetData/,
    FontIndex=3,
    AutoFakeBold=2
  ]{PingFang.ttc}
}
\usetikzlibrary{arrows.meta,bending,positioning,calc,backgrounds,fit}

% Low-saturation academic palette.
\definecolor{ink}{HTML}{27364A}
\definecolor{muted}{HTML}{748092}
\definecolor{mainLine}{HTML}{4F667C}
\definecolor{mainFill}{HTML}{EAF0F4}
\definecolor{heroFill}{HTML}{F1F3F6}
\definecolor{heroBorder}{HTML}{3F566C}
\definecolor{blueLine}{HTML}{6D879E}
\definecolor{blueFill}{HTML}{E7EEF4}
\definecolor{purpleLine}{HTML}{817594}
\definecolor{purpleFill}{HTML}{EEEAF2}
\definecolor{goldLine}{HTML}{9B875B}
\definecolor{goldFill}{HTML}{F2EEE3}
\definecolor{tealLine}{HTML}{648D8A}
\definecolor{tealFill}{HTML}{E5F0EE}
\definecolor{roseLine}{HTML}{9A777B}
\definecolor{roseFill}{HTML}{F3E9EA}
\definecolor{greenLine}{HTML}{71896E}
\definecolor{greenFill}{HTML}{EAF0E8}
\definecolor{feedback}{HTML}{7C6D8D}
\definecolor{handoff}{HTML}{8A777A}

\pgfdeclarelayer{background}
\pgfsetlayers{background,main}

\tikzset{
  every node/.style={
    font=\sffamily\fontsize{8.8}{10.5}\selectfont,
    text=ink
  },
  herozone/.style={
    rectangle,
    rounded corners=6pt,
    draw=heroBorder,
    fill=heroFill,
    line width=1.4pt,
    minimum width=14.2cm,
    minimum height=8.12cm,
    inner sep=0pt
  },
  panel/.style={
    rectangle,
    rounded corners=5pt,
    line width=0.7pt,
    minimum width=5.90cm,
    minimum height=2.40cm,
    text width=5.18cm,
    align=left,
    inner sep=7pt,
    font=\sffamily\fontsize{8.4}{10.6}\selectfont
  },
  summarybar/.style={
    rectangle,
    rounded corners=4pt,
    draw=muted!55,
    fill=mainFill!55,
    line width=0.6pt,
    minimum width=27.20cm,
    minimum height=1.15cm,
    align=center,
    inner sep=6pt,
    font=\sffamily\fontsize{9.0}{10.8}\selectfont\bfseries
  },
  flowbox/.style={
    rectangle,
    rounded corners=4pt,
    draw=mainLine,
    fill=white,
    line width=0.7pt,
    minimum width=2.50cm,
    minimum height=1.15cm,
    text width=2.08cm,
    align=center,
    inner sep=3.5pt
  },
  lowerbox/.style={
    flowbox,
    minimum width=2.50cm,
    text width=2.08cm
  },
  % Canonical arrow styles adapted directly from tikz-template.tex.
  arrow/.style={
    -{Stealth[length=5pt 1.5,width'=0pt 0.6,sep=0pt 0.5]},
    line width=1.0pt,
    shorten >=2pt,
    shorten <=1pt,
    color=mainLine
  },
  arrow thick/.style={arrow,line width=1.6pt},
  arrow thin/.style={arrow,line width=0.6pt},
  arrow short/.style={
    -{Stealth[length=3pt 1.0,width'=0pt 0.5,sep=0pt 0.3]},
    line width=0.8pt,
    shorten >=1pt,
    shorten <=0pt,
    color=mainLine
  },
  residual/.style={
    -{Stealth[length=4pt 1.2,width'=0pt 0.6]},
    line width=0.9pt,
    densely dashed,
    shorten >=2pt,
    color=feedback
  },
  residual line/.style={
    line width=0.9pt,
    densely dashed,
    color=feedback
  },
  residual straight/.style={residual,rounded corners=0pt},
  leader/.style={
    densely dotted,
    line width=0.32pt,
    color=muted!55,
    shorten >=2pt,
    shorten <=2pt
  },
  fan_stub/.style={
    -{Stealth[length=5pt 1.5,width'=0pt 0.6,sep=0pt 0.5]},
    line width=1.0pt,
    shorten >=2pt,
    shorten <=0pt,
    color=mainLine
  }
}

\begin{document}
\sffamily
\begin{tikzpicture}

% Surrounding chapter panels: three left, three right, all placed by construction.
\node[panel,anchor=north west,fill=blueFill,draw=blueLine] (panel1) at (0,0)
  {{\fontsize{9.1}{10.6}\selectfont\bfseries §4 任务·环境·接口}\\[3pt]
   Web·Mobile·Desktop；\\
   GUI·CLI·API·MCP；Hybrid 路由\\[3pt]
   {\color{muted}\fontsize{8.0}{9.2}\selectfont RQ2·RQ3·RQ8}};
\node[panel,below=0.40cm of panel1,fill=tealFill,draw=tealLine] (panel5)
  {{\fontsize{9.1}{10.6}\selectfont\bfseries §8 评测}\\[3pt]
   grounding→轨迹→在线 benchmark；\\
   metrics·verifier·可复现性\\[3pt]
   {\color{muted}\fontsize{8.0}{9.2}\selectfont RQ9}};
\node[panel,below=0.40cm of panel5,fill=goldFill,draw=goldLine] (panel2)
  {{\fontsize{9.1}{10.6}\selectfont\bfseries §5 数据}\\[3pt]
   \mbox{理解→grounding→轨迹→推理} traces；\\
   失败·偏好·个性化\\[3pt]
   {\color{muted}\fontsize{8.0}{9.2}\selectfont RQ7}};

\node[herozone,right=0.60cm of panel1.north east,anchor=north west] (hero) {};

\node[panel,right=0.60cm of hero.north east,anchor=north west,
      fill=purpleFill,draw=purpleLine] (panel3)
  {{\fontsize{9.1}{10.6}\selectfont\bfseries §6 模型与系统}\\[3pt]
   Native 端到端·模块化·Multi-agent·Hybrid；\\
   感知·规划·记忆·验证·安全\\[3pt]
   {\color{muted}\fontsize{8.0}{9.2}\selectfont RQ4·RQ5}};
\node[panel,below=0.40cm of panel3,fill=roseFill,draw=roseLine] (panel6)
  {{\fontsize{9.1}{10.6}\selectfont\bfseries §9–10 部署与开放问题}\\[3pt]
   产业格局·治理；十大开放方向\\[3pt]
   {\color{muted}\fontsize{8.0}{9.2}\selectfont RQ10·RQ11·RQ12}};
\node[panel,below=0.40cm of panel6,fill=greenFill,draw=greenLine] (panel4)
  {{\fontsize{9.1}{10.6}\selectfont\bfseries §7 学习与优化}\\[3pt]
   Pre-train→SFT→RL（online·RLVR）→\\
   自我进化\\[3pt]
   {\color{muted}\fontsize{8.0}{9.2}\selectfont RQ6}};

% Small global title at the upper-left.
\node[
  anchor=south west,
  font=\sffamily\fontsize{10.0}{12.0}\selectfont\bfseries,
  text=heroBorder
] at ([yshift=0.63cm]panel1.north west)
  {Computer-Use Agents: A Unified Survey — 全文组织};

% Full-width summary bar.
\node[summarybar,below=0.45cm of panel2.south west,anchor=north west] (summary)
  {Accountable-state thesis：来源可追溯 · 结果可核验 · 失败可恢复的状态转移是 CUA 的优化单元 (§1、§11)};

% Central hero.
\node[
  anchor=south,
  font=\sffamily\fontsize{10.2}{12.0}\selectfont\bfseries,
  text=heroBorder
] (heroTitle) at ([yshift=0.12cm]hero.north) {CUA 执行闭环 (§3)};

% Main loop, first row.
\coordinate (mainStart) at ([xshift=0.725cm,yshift=-1.35cm]hero.north west);
\node[flowbox,anchor=north west,fill=blueFill,draw=blueLine] (intent) at (mainStart)
  {用户意图\\I};
\node[flowbox,right=0.85cm of intent,fill=blueFill,draw=blueLine] (obs)
  {Observation / Belief\\O};
\node[flowbox,right=0.85cm of obs,fill=purpleFill,draw=purpleLine] (state)
  {Explicit Task State\\S};
\node[flowbox,right=0.85cm of state,fill=goldFill,draw=goldLine] (plan)
  {Planning / Policy\\P};

% Main loop, second row: semantic flow runs right-to-left.
\node[flowbox,below=1.35cm of intent,fill=roseFill,draw=roseLine] (recovery)
  {Recovery / Abstention\\R};
\node[flowbox,right=0.85cm of recovery,fill=tealFill,draw=tealLine] (verifier)
  {Verifier\\V};
\node[flowbox,right=0.85cm of verifier,fill=mainFill,draw=mainLine] (env)
  {Environment\\E};
\node[flowbox,right=0.85cm of env,fill=roseFill,draw=roseLine] (action)
  {Semantic GUI/API\\Action A};

% Learning / factory / handoff layer.
\node[lowerbox,below=0.90cm of recovery,fill=goldFill,draw=goldLine] (factory)
  {Data / Task Factory\\D};
\node[lowerbox,right=2.50cm of factory,fill=greenFill,draw=greenLine] (learning)
  {Learning\\L};
\node[lowerbox,right=2.50cm of learning,fill=roseFill,draw=handoff] (human)
  {Human Handoff\\H};

% Chapter leaders are drawn before semantic arrows, so primary flow remains on top.
% §4 -> O and A: shared trunk at the hero's left margin.
\coordinate (panel1Trunk) at ([xshift=0.22cm]hero.west |- panel1.east);
\coordinate (panel1ObsRail) at ([xshift=0.22cm,yshift=-0.95cm]hero.north west);
\coordinate (panel1ActionRail) at ([xshift=0.22cm,yshift=-3.05cm]hero.north west);
\coordinate (panel1ObsTop) at (obs.north |- panel1ObsRail);
\coordinate (panel1ActionTarget) at ([xshift=-0.45cm]action.north);
\draw[leader] (panel1.east) -- (panel1Trunk);
\draw[leader] (panel1ObsRail) -- (panel1ActionRail);
\draw[leader] (panel1ObsRail) -- (panel1ObsTop) -- (obs.north);
\draw[leader]
  (panel1ActionRail) -- (panel1ActionTarget |- panel1ActionRail)
  -- (panel1ActionTarget);

% §5 -> D; §6 -> P; §7 -> L; §8 -> V.
\draw[leader] (panel2.east) -- (panel2.east |- factory.west) -- (factory.west);
\coordinate (panel3TopRail) at ([yshift=0.22cm]hero.north);
\coordinate (panel3PlanTarget) at ([xshift=0.55cm]plan.north);
\draw[leader]
  (panel3.north west) -- (panel3.north west |- panel3TopRail)
  -- (panel3PlanTarget |- panel3TopRail) -- (panel3PlanTarget);
\coordinate (panel4BottomRail) at ([yshift=-0.18cm]hero.south);
\coordinate (panel4LearnTarget) at ([xshift=0.50cm]learning.south);
\draw[leader]
  (panel4.south west) -- (panel4.south west |- panel4BottomRail)
  -- (panel4LearnTarget |- panel4BottomRail) -- (panel4LearnTarget);
\coordinate (panel5Entry) at ([xshift=-0.30cm]hero.west |- panel5.east);
\coordinate (panel5VerifierRail) at ([yshift=-3.15cm]hero.north);
\draw[leader]
  (panel5.east) -- (panel5Entry)
  -- (panel5Entry |- panel5VerifierRail)
  -- (verifier.north |- panel5VerifierRail) -- (verifier.north);

% §9--10 -> R and H: shared lower spine with two short target stubs.
\coordinate (panel6SpineRight) at ([xshift=0.25cm,yshift=-5.35cm]hero.north east);
\coordinate (panel6SpineLeft) at (recovery.south |- panel6SpineRight);
\coordinate (panel6SpineHuman) at (human.north |- panel6SpineRight);
\draw[leader]
  (panel6.south west) -- (panel6.south west |- panel6SpineRight)
  -- (panel6SpineRight);
\draw[leader] (panel6SpineLeft) -- (panel6SpineRight);
\draw[leader] (panel6SpineLeft) -- (recovery.south);
\draw[leader] (panel6SpineHuman) -- (human.north);

% Forward path: all adjacent links are <1.5 cm and therefore use the small tip.
\draw[arrow short] (intent.east) -- (obs.west);
\draw[arrow short] (obs.east) -- (state.west);
\draw[arrow short] (state.east) -- (plan.west);
\draw[arrow] (plan.south) -- (action.north);
\draw[arrow short] (action.west) -- (env.east);
\draw[arrow short] (env.west) -- (verifier.east);
\draw[arrow short] (verifier.west) -- (recovery.east);

% R -> P feedback uses the lower of two dedicated rails between the loop rows.
\coordinate (feedbackRailLeft) at ([yshift=0.34cm]recovery.north);
\coordinate (feedbackPlanTarget) at ([xshift=-0.56cm]plan.south);
\coordinate (feedbackRailRight) at (feedbackPlanTarget |- feedbackRailLeft);
\draw[residual line]
  (recovery.north) -- (recovery.north |- feedbackRailLeft) -- (feedbackRailRight);
\draw[residual straight] (feedbackRailRight) -- (feedbackPlanTarget);

% R -- H (dashed) crosses the lower corridor to a gap before H.
\coordinate (handoffShelf) at ([yshift=-0.40cm]recovery.south);
\coordinate (handoffGap) at ($(learning.east)!0.5!(human.west)$);
\coordinate (handoffGapShelf) at (handoffGap |- handoffShelf);
\coordinate (handoffBottom) at ([yshift=0.27cm]hero.south);
\coordinate (handoffGapBottom) at (handoffGap |- handoffBottom);
\coordinate (handoffHumanBottom) at (human.south |- handoffBottom);
\draw[residual line,color=handoff]
  (recovery.south) -- (recovery.south |- handoffShelf)
  -- (handoffGapShelf) -- (handoffGapBottom) -- (handoffHumanBottom);
\draw[residual straight,color=handoff] (handoffHumanBottom) -- (human.south);

% H -> S uses a dedicated right/top rail, avoiding the A -> E horizontal path.
\coordinate (humanRightRail) at ([xshift=-0.50cm]hero.east);
\coordinate (humanTopRight) at ([xshift=-0.50cm,yshift=-0.95cm]hero.north east);
\coordinate (humanStateTop) at (state.north |- humanTopRight);
\draw[line width=0.6pt,color=mainLine]
  (human.east) -- (humanRightRail |- human.east) -- (humanTopRight);
\draw[line width=0.6pt,color=mainLine] (humanTopRight) -- (humanStateTop);
\draw[arrow short] (humanStateTop) -- (state.north);

% D -> L and V -> L.
\draw[arrow,color=greenLine] (factory.east) -- (learning.west);
\coordinate (learnRail) at ([yshift=-0.64cm]verifier.south);
\draw[line width=2.8pt,color=heroFill]
  (verifier.south) -- (verifier.south |- learnRail)
  -- (learning.north |- learnRail) -- (learning.north);
\draw[arrow short,color=greenLine]
  (verifier.south) -- (verifier.south |- learnRail)
  -- (learning.north |- learnRail) -- (learning.north);

% L -> O/P rises through the aligned E/A and S/P gap. A narrow hero-fill
% underpass makes each unavoidable orthogonal crossing visibly non-junctional.
\coordinate (learnGap) at ($(env.east)!0.5!(action.west)$);
\coordinate (learnBottom) at ([yshift=0.52cm]hero.south);
\coordinate (learnFromL) at (learning.south |- learnBottom);
\coordinate (learnGapBottom) at (learnGap |- learnBottom);
\coordinate (learnRailLeft) at ([yshift=0.78cm]recovery.north);
\coordinate (learnObsTarget) at ([xshift=-0.22cm]obs.south);
\coordinate (learnPlanTarget) at ([xshift=-0.18cm]plan.south);
\coordinate (learnRailJunction) at (learnGap |- learnRailLeft);
\coordinate (learnRailObs) at (learnObsTarget |- learnRailLeft);
\coordinate (learnRailPlan) at (learnPlanTarget |- learnRailLeft);
\draw[line width=2.8pt,color=heroFill]
  (learning.south) -- (learnFromL) -- (learnGapBottom)
  -- (learnRailJunction) -- (learnRailPlan);
\draw[line width=2.8pt,color=heroFill] (learnRailJunction) -- (learnRailObs);
\draw[line width=0.6pt,color=greenLine]
  (learning.south) -- (learnFromL) -- (learnGapBottom)
  -- (learnRailJunction) -- (learnRailPlan);
\draw[line width=0.6pt,color=greenLine] (learnRailJunction) -- (learnRailObs);
\draw[arrow short,color=greenLine] (learnRailObs) -- (learnObsTarget);
\draw[arrow short,color=greenLine] (learnRailPlan) -- (learnPlanTarget);

\end{tikzpicture}
\end{document}
